\documentclass[twocolumn]{aastex631}
\begin{document}
\title{The reionization ionizing-photon-budget shortfall for star-forming galaxies is not robust to systematics at z\~{}10 (optimistic systematic corner)}
\author{NebulaMind Lab (autonomous pipeline)}
\affiliation{NebulaMind Lab, \url{https://lab.nebulamind.net}}
\begin{abstract}
Reionization ionizing-photon-budget reconciliation at z\~{}10: star-forming galaxies require f\_esc=0.111 (+0.096/-0.051) to close the budget (Madau-Dickinson SFRD, log xi\_ion=25.5+/-0.15, clumping C=2-5, JWST-SFRD tail), versus indirect-proxy-inferred f\_esc=0.080 (+0.147/-0.051) from LzLCS O32/beta calibrations. Median delta(required-inferred)=+0.026 dex-frac (16-84\%: -0.114 to +0.130); 61\% of the systematic MC shows a shortfall. Conclusion: the budget CLOSES within the systematic, and the sign holds under both O32 and beta calibrations. The result is bounded by the xi\_ion x clumping x proxy-calibration systematic, not by statistics. Generated autonomously from public data (jwst) via the NebulaMind Lab runner: a bounded, reproducible, descriptive study.
\end{abstract}
\section{Introduction}
Recent studies have highlighted potential discrepancies in the ionizing photon budget during the reionization epoch, suggesting that current observations may not account for the necessary photons to drive this process [Muñoz2024]. This has led to concerns about a "photon budget crisis" and the need for further investigation into the sources of these photons. Previous work has explored various aspects of reionization, including the role of galaxies in powering this process [Duncan2015] and the development of models to better understand the timing and mechanisms involved [Park2022]. However, a comprehensive analysis of the ionizing photon budget is still required to reconcile these findings.
\section{Data and method}
To address this issue, we employ a literature-anchored budget calculation that does not rely on survey catalog data. Instead, we utilize the cosmic star formation rate density (SFRD) from Madau \& Dickinson (2014), along with published values for xi\_ion and O32/beta f\_esc proxy calibrations [LzLCS: Chisholm+22, Flury+22; Simmonds+24]. By adopting these established values, we aim to systematically reconcile the ionizing photon budget during reionization.
\section{Result}
Our analysis reveals that star-forming galaxies must have an escape fraction of f\_esc=0.111 (+0.096/-0.051) at z\~{}10 to close the ionizing photon budget, assuming a Madau-Dickinson SFRD, log xi\_ion=25.5±0.15, and clumping factor C=2-5. This result is compared to indirect-proxy-inferred f\_esc values of 0.080 (+0.147/-0.051) derived from LzLCS O32/beta calibrations. The median difference between the required and inferred escape fractions is +0.026 dex-frac, with a range of -0.114 to +0.130 (16-84\% confidence interval). Notably, 61\% of our systematic Monte Carlo simulations indicate a shortfall in the photon budget.
\begin{figure}\centering\includegraphics[width=\columnwidth]{result.png}\caption{The reionization ionizing-photon-budget shortfall for star-forming galaxies is not robust to systematics at z\~{}10 (optimistic systematic corner)}\end{figure}
\section{Caveats}
It is essential to acknowledge the limitations of this study. Our analysis relies on automated, single-selection, and uncalibrated measurements, which may introduce biases or inaccuracies. The escape fraction values derived from O32/beta calibrations are subject to uncertainties in the underlying assumptions and models used to establish these proxies. Furthermore, our reconciliation of the ionizing photon budget is sensitive to the choice of xi\_ion and clumping factor, highlighting the need for more precise constraints on these parameters. Despite these caveats, our findings provide valuable insights into the reionization process and underscore the importance of continued research in this area.
\section*{References}
\footnotesize
[Muoz2024] Muñoz et al. (2024). Reionization after JWST: a photon budget crisis?. bibcode 2024MNRAS.535L..37M \par [Duncan2015] Duncan et al. (2015). Powering reionization: assessing the galaxy ionizing photon budget at z < 10. bibcode 2015MNRAS.451.2030D \par [Park2022] Park et al. (2022). Calibrating excursion set reionization models to approximately conserve ionizing photons. bibcode 2022MNRAS.517..192P \par [Davies2021] Davies et al. (2021). The Predicament of Absorption-dominated Reionization: Increased Demands on Ionizing Sources. bibcode 2021ApJ...918L..35D \par [Madau2017] Madau (2017). Cosmic Reionization after Planck and before JWST: An Analytic Approach. bibcode 2017ApJ...851...50M

\end{document}
